% Joint root/child research derivation for R75. Not a sealed source edit.
\section{An unbounded nonabelian family with a strict positivity cost}
\label{sec:r75-symmetric-return}

We apply the multiplicity theory to all symmetric groups $S_n$ with
$n\ge5$. The natural permutation character $\chi_n(g)$ counts the
fixed points of $g$ on $[n]$. Consider the binary response whose first
coordinate is
\begin{equation}\label{eq:r75-symmetric-return}
 \beta_t(g)=\frac{1-t}{n}+\frac{t}{n}\chi_n(g),
 \qquad b_t=(\beta_t,1-\beta_t),\qquad0\le t\le1.
\end{equation}
Its group average is $1/n$, and at the identity it equals
$[1+(n-1)t]/n$. This supplies a family with growing group, observable
rank, and exact memory, rather than a fixed-dimensional illustration.

\begin{theorem}[Sharp nonabelian endpoint memory]
\label{thm:r75-symmetric-return}
For $n\ge5$ and $t>0$, the response \eqref{eq:r75-symmetric-return}
has observable group $S_n$, residual cardinality $n!$, and observable
dimension $1+(n-1)^2$. Its permutation-module bound is
\[
 \nu_{S_n}(V_{b_t})=n(n-1).
\]
The unrestricted positive core and pure-start memory satisfy
\[
 \begin{array}{c|c|c}
 \text{contrast}&\kappa_{S_n}(b_t)&m(b_t)\\\hline
 0<t\le1/(n-1)&n(n-1)&n(n-1)+1\\
 t=1&n^2&n^2+1.
 \end{array}
\]
At $t=0$ both minima are one. Moreover there is a number
$\tau_n<1$ such that $\kappa_{S_n}(b_t)=n^2$ for every
$\tau_n<t\le1$. No formula for the intermediate transition points
is asserted here.
\end{theorem}
\begin{proof}
The real natural permutation module is the direct sum of the
invariant line and the $(n-1)$-dimensional standard representation
$\mathrm{Std}_n$. Its complexification is irreducible on that
zero-sum subspace. One direct proof uses the transpositions: from
a nonconstant vector in an invariant subspace, a vector
$v-(ij)v$ produces a nonzero multiple of $e_i-e_j$, whose orbit
spans the whole zero-sum space. The standard representation is
faithful, because a permutation fixing every difference $e_i-e_j$
fixes all points. Matrix-coefficient orthogonality gives the
nontrivial block
\[
 \widehat\beta_t(\mathrm{Std}_n)
   =\frac{t}{n(n-1)}I_{n-1},
\]
and all other nontrivial irreducible blocks vanish. Therefore the
observable module is
$\boldsymbol1\oplus\mathrm{Std}_n^{\oplus(n-1)}$, of dimension
$1+(n-1)^2$. Equality of two residuals forces equality of their
standard representation matrices because this Fourier block is
invertible; faithfulness then makes the group elements equal.
This proves the claims about the group and residual count.

We next record the subgroup index estimate needed for an unrestricted
memory converse. For $K\le S_n$, let $r$ be the number of its orbits
on $[n]$. The multiplicity of $\mathrm{Std}_n$ in
$\mathbb C[K\backslash S_n]$ is
\[
 m(K)=\dim\mathrm{Std}_n^K=r-1.
\]
If $m(K)>0$, then
\begin{equation}\label{eq:r75-index-estimate}
 [S_n:K]\ge n\,m(K).
\end{equation}
More sharply, unless $K$ is a full point stabilizer,
\begin{equation}\label{eq:r75-index-extra}
 [S_n:K]\ge n\{m(K)+1\}.
\end{equation}
To prove these, let the natural orbit sizes of $K$ be
$a_1,\ldots,a_r$. Since $K$ is contained in the product of the
symmetric groups on those orbits,
\[
 [S_n:K]\ge\frac{n!}{a_1!\cdots a_r!}.
\]
For $r\ge3$, this is at least $n!/(n-r+1)!$: transferring one
unit from the smaller of two nonunit parts to the larger increases
their factorial product until only one nonunit part remains. The
last lower bound is at least $nr$ for $n\ge5$. Indeed at $r=3$
it is $n(n-1)\ge3n$, and increasing $r$ multiplies the bound by
at least two while $r$ increases by one. If $r=2$ and one part
is a singleton, $K$ lies in a point stabilizer; it is that full
stabilizer of index $n$, or has index at least $2n$. If both
parts have size at least two, the index is at least
$\binom n2\ge2n$. These cases prove both inequalities.

A $G$-set containing $V_{b_t}$ must have total standard multiplicity
at least $n-1$. Summing \eqref{eq:r75-index-estimate} over its
standard-carrying orbits gives at least $n(n-1)$ points.
Exactly $n-1$ natural orbits supply this multiplicity, so the module
minimum is $n(n-1)$.

There is a positive realization attaining this bound on an explicit
contrast interval. Use $n-1$ natural orbits indexed by
$i=1,\ldots,n-1$, each with initial mass $1/(n-1)$ concentrated
at its point $i$. Give orbit $i$ the scalar decoder
\begin{equation}\label{eq:r75-symmetric-decoder}
 D_i(j)=\frac1n+\frac{(n-1)t}{n}
                 \{\mathbf1_{\{j=i\}}-\mathbf1_{\{j=n\}}\}.
\end{equation}
For $0\le t\le1/(n-1)$ these numbers lie in $[0,1]$.
Since a permutation has exactly one preimage of $n$,
\[
 \frac1{n-1}\sum_{i=1}^{n-1}D_i(g(i))
 =\frac1n+\frac tn
   \sum_{i=1}^{n-1}
    \{\mathbf1_{\{g(i)=i\}}-\mathbf1_{\{g(i)=n\}}\}
 =\frac1n+\frac tn\{\chi_n(g)-1\}=\beta_t(g).
\]
The module bound proves optimality on this interval.

At full contrast, $n$ natural orbits realize $\beta_1$: give
each mass $1/n$, start orbit $i$ at $i$, and use the indicator
decoder of $i$. Hence $\kappa_{S_n}(b_1)\le n^2$.
Suppose, to the contrary, that a factorization exists on fewer
than $n^2$ points. The total standard multiplicity is at least
$n-1$. If any standard-carrying orbit were not natural,
\eqref{eq:r75-index-extra} and the ordinary estimates on the
others would force at least $n(n-1)+n=n^2$ points. Thus all
standard-carrying orbits are natural. There must be exactly
$n-1$ of them: each supplies multiplicity one and costs $n$
points. Every other orbit has total size less than $n$.

Write the factorization as a mixture over its positive-mass
transitive orbits. Since $\beta_1(e)=1$, each component's
decoder equals one on its initial support. On a transitive orbit
of degree $d$ its mean is therefore at least $1/d$.
The natural components have mean at least $1/n$; any additional
orbit here has $d<n$ and would have mean strictly greater than
$1/n$. The mean of $\beta_1$ is exactly $1/n$, so all additional
orbits have zero initial mass and every active natural component
has mean $1/n$. Lemma~\ref{lem:r75-return-rigidity} now makes
each component a deterministic return law. Let $w_i$ be the
aggregate mass of the components starting at natural point $i$.
Evaluating at a transposition $(ij)$ gives
\[
 1-w_i-w_j=\frac{n-2}{n},
 \qquad w_i+w_j=\frac2n.
\]
For $n\ge3$ these equations imply $w_i=1/n$ for every $i$.
Thus at least $n$ positive components are needed, contradicting
the availability of only $n-1$ natural orbits. This proves
$\kappa_{S_n}(b_1)=n^2$.

The full-contrast construction also realizes every $b_t$ by
mixing its decoder with the constant $1/n$, so $\kappa\le n^2$
at every contrast. The attained-threshold theorem applied to
budget $n^2-1$ gives a closed feasible interval
$[0,\tau_n]$. Full contrast is not in it, hence $\tau_n<1$;
above it the matching value is $n^2$. The pure-start formula
$m=\min\{n!,1+\kappa\}$ gives the displayed values since
$n^2+1<n!$ for $n\ge5$.
\end{proof}

\paragraph{Prepared quantum realization.}
The family has a finite-dimensional prepared unitary realization,
so it remains inside the experiment framework. On
$\mathbb C^n\otimes\mathbb C^n$ take
\[
 \rho=\frac1n\sum_{i=1}^n|ii\rangle\langle ii|,
 \qquad U_g=P_g\otimes I,
 \qquad E=\sum_{i=1}^n|ii\rangle\langle ii|.
\]
Then $\operatorname{tr}(E U_g\rho U_g^*)=\chi_n(g)/n$.
The binary effect $E_t=tE+(1-t)I/n$ produces
\eqref{eq:r75-symmetric-return}. Any finite generating command set
for $S_n$ supplies one common experiment valid for all words.
The pure maximally entangled preparation
$n^{-1/2}\sum_i|ii\rangle$ gives the same terminal probabilities
with this effect, if a pure preparation is desired. The statement
is a physical realization of the response, not a claim that this
Hilbert-space dimension is minimal or that it establishes a
quantum-versus-classical separation under a common resource metric.

The gap $n$ between the positive core at full contrast and its
module minimum is unbounded. The relative ratio $n/(n-1)$ tends
to one; no growing relative separation is asserted. The explicit
low-contrast interval is a proved attainment interval, not a claim
that $1/(n-1)$ is the exact first transition.
